\subsubsection{Sobolev 嵌入定理}\label{subsec:sobolev-imbedding-theorem}

先定义术语 ``嵌入''.

\begin{Definition}
\label{def:continuous-imbedding}
设 $X$ 和 $Y$ 为 Banach 空间, 若 $X\subset Y$ 且存在常数 $C$, 使
\begin{equation}\label{eq:continuous-embedding-norm-bound}
\|u\|_Y\leq C\|u\|_X\quad\forall u\in X.
\end{equation}
则称 $X$ \textbf{连续嵌入} $Y$, 记为\[X\hookrightarrow Y\]
\end{Definition}

\begin{Remark}
\label{rem:continuous-imbedding-bounded-identity}
若 $X\subset Y$, 则恒等算子 $I:X\to Y$ 有良好定义. 条件 \eqref{eq:continuous-embedding-norm-bound} 表示, 算子的输出由输入的范数乘以某个常数控制. 因而恒等算子是有界算子.
\end{Remark}

Sobolev 嵌入定理表明, 若 $\Omega$ 足够好, 则对充分大的 $k$, $H^k(\Omega)$ 嵌入某个有界连续函数空间. 先考察 $\Omega=\mathbb{R}^m$ 的情形. 作为准备, 陈述以下引理.

\begin{Lemma}
\label{lem:integrable-fourier-transform-continuity}
若 $\mathcal{F}[u]\in L^1(\mathbb{R}^m)$, 则 $u$ 为连续有界函数. 此外, $\|u\|_{\infty}\leq(2\pi)^{-m/2}\|\mathcal{F}[u]\|_1$.
\end{Lemma}

\begin{Proof}
若 $\mathcal{F}[u]\in\mathcal{S}(\mathbb{R}^m)$, 则 $u\in\mathcal{S}(\mathbb{R}^m)$, 且
\begin{equation}\label{eq:integrable-fourier-transform-uniform-bound}
\begin{aligned}
\|u\|_{\infty}&=\sup_{\bm{x}\in\mathbb{R}^m}(2\pi)^{-m/2}\left|\int_{\mathbb{R}^m}\exp(i\bm{\xi}\cdot\bm{x})\mathcal{F}[u](\bm{\xi})\,\mathrm{d}\bm{\xi}\right|\\
&\leq(2\pi)^{-m/2}\|\mathcal{F}[u]\|_1.
\end{aligned}
\end{equation}
由于 $\mathcal{S}(\mathbb{R}^m)$ 在 $L^1(\mathbb{R}^m)$ 中稠密, 取极限即得引理.
\end{Proof}

\begin{Theorem}{Sobolev 嵌入定理}
\label{thm:sobolev-full-space-continuous-embedding}
设 $s>m/2$. 那么\[H^s(\mathbb{R}^m)\hookrightarrow C_b(\mathbb{R}^m).\]即\[H^s(\mathbb{R}^m)\subset C_b(\mathbb{R}^m)\]且存在常数 $C$, 使对每个 $u\in H^s(\mathbb{R}^m)$ 有 $\|u\|_{\infty}\leq C\|u\|_{s,2}$.
\end{Theorem}

\begin{Proof}
设 $u\in H^s(\mathbb{R}^m)$. 那么 $\mathcal{F}[u]\in L_s^2(\mathbb{R}^m)$. 因而由前一引理, 只需证明 $L_s^2(\mathbb{R}^m)$ 连续嵌入 $L^1(\mathbb{R}^m)$. 令 $w(\bm{\xi})=1+|\bm{\xi}|^{2s}$. 那么由 Hölder 不等式得到
\begin{equation}\label{eq:sobolev-embedding-fourier-holder-weight}
\|\mathcal{F}[u]\|_1\leq\|w^{1/2}\mathcal{F}[u]\|_2\|w^{-1/2}\|_2.
\end{equation}
由于已假设 $s>m/2$, $w^{-1/2}$ 属于 $L^2(\mathbb{R}^m)$.
\end{Proof}

\begin{Corollary}
\label{cor:sobolev-extension-continuity}
假设 $k>m/2$, 且 $\Omega$ 具有 $k$ 阶延拓性质. 那么\[H^k(\Omega)\hookrightarrow C_b(\overline\Omega).\]若 $k>(m/2)+j$, 则\[H^k(\Omega)\hookrightarrow C_b^j(\overline\Omega).\]
\end{Corollary}

\begin{Remark}
\label{rem:sobolev-embedding-optimal-threshold}
在 Theorem~\ref{thm:sobolev-full-space-continuous-embedding} 中, 界 $s>m/2$ 是最优的; 参见 Problem~\ref{ex:sobolev-embedding-critical-counterexample}.
\end{Remark}

\Needspace{4\baselineskip}
\begin{Example}
\label{ex:sobolev-jump-discontinuity}
在 Example~\ref{ex:heaviside-absolute-value-sobolev-membership} 中注意到, Heaviside 函数 $H(x)$ 不属于 $H^1(-1,1)$. 同样的论证可用于证明, 任何具有跳跃间断点的函数都不属于 $H^1(-1,1)$. Sobolev 嵌入定理给出强得多的结果: 若 $s>1/2$, 则 $H^s(-1,1)$ 中每个函数都有界且连续.
\end{Example}

\Needspace{4\baselineskip}
\begin{Example}
\label{ex:sobolev-radial-power-embedding}
设 $B$ 为 $\mathbb{R}^m$ 中半径为一的球. 在 Example~\ref{ex:radial-power-sobolev-membership} 中证明了, $|\bm{x}|^\alpha\in W^{k,p}(B)$ 当且仅当 $\alpha>k-m/p$. 若令 $p=2$, 则注意到 Sobolev 嵌入定理所要求的条件 $k>m/p$, 恰好保证 $\alpha>0$, 从而这里的径向函数连续.
\end{Example}

\Needspace{4\baselineskip}
\begin{Example}
\label{ex:sobolev-three-dimensional-critical-order}
在 $\mathbb{R}^3$ 中, 临界 ``导数阶数'' 为 $s=3/2$. $H^1(\mathbb{R}^3)$ 中的函数未必连续, 但 $H^2(\mathbb{R}^3)$ 中的所有函数都有界且连续.
\end{Example}

\begin{Remark}
\label{rem:sobolev-general-p-embeddings}
可以建立 $W^{k,p}$ 空间的更一般嵌入定理. 特别地, 可以证明\[W^{k,p}(\mathbb{R}^m)\hookrightarrow L^{mp/(m-kp)}(\mathbb{R}^m)\quad\text{for }kp<m,\]以及\[W^{k,p}(\mathbb{R}^m)\hookrightarrow C_b(\mathbb{R}^m)\quad\text{for }kp>m,\]且这些嵌入中的每一个都连续. 这同样可推广到更一般的区域 $\Omega$, 例如具有延拓性质的区域.
\end{Remark}
